% Part 1: course order, 8 October through 9 October 07:59.
\section{从无线网络到一条无线链路}
\subsection{先看应用需要什么，再看网络怎样组成}
移动与传感器网络（Mobile and Sensor Networks）的起点不是某个通信公式，而是设备为什么需要连接。环境传感器周期性报告少量测量值，视频设备连续上传大量数据，车辆则需要在移动过程中交换信息。它们都使用无线网络（Wireless Network），却不要求相同的服务。视频重视吞吐量（Throughput），也就是单位时间内实际交付的数据量；及时控制重视时延（Latency），也就是信息从产生或发送到到达所经历的时间；电池供电的传感器还重视能耗（Energy Consumption）。可靠性（Reliability）关注信息能否按要求成功交付，例如用一定条件下成功交付的比例来衡量。设备数量、移动速度和覆盖范围，也会改变网络的设计。

因此，“通信速率越高，网络就越好”并不是普遍成立的判断。一台每天只上报少量数据的传感器，可能更需要几年不用换电池；一台高速移动的终端，需要在离开原来覆盖区以后继续保持业务。课程先介绍这些不同需求，是为了说明后面的调制、传播、错误控制和接入机制，都是解决具体约束的工具。

一个无线主机（Wireless Host）可以是手机、笔记本，也可以是没有屏幕的传感器。主机通过无线链路（Wireless Link）与另一个节点交换数据。基础设施网络（Infrastructure Network）中，终端通常连接到接入点（Access Point, AP）或基站（Base Station, BS），再由它们连接其他网络。接入点是无线局域网中常见的连接节点，基站是蜂窝网络中的常见连接节点，两者在这里承担相似的接入作用，但不意味着它们的协议和内部结构相同。

自组织网络（Ad Hoc Network）则允许设备直接通信，不要求每次交换都经过固定接入点。混合架构（Hybrid Architecture）把直接通信与基础设施接入结合起来，例如车辆之间直接交换消息，同时通过路侧设备进入更大的网络。网状网络（Mesh Network）进一步让节点相互转发数据：不是每个节点都有独立互联网出口，数据可以经过其他节点，到达具有出口的网关（Gateway）。这种架构把“能与邻居通信”和“能直接连接互联网”分成了两个问题。

\begin{figure}[htbp]
\centering
\begin{tikzpicture}[>=stealth, every node/.style={font=\small}]
\node[draw,rounded corners] (a) at (0,0) {终端};
\node[draw,rounded corners] (b) at (2.4,0) {接入点};
\node[draw,rounded corners] (c) at (5,0) {其他网络};
\draw[<->] (a)--(b); \draw[<->] (b)--(c);
\node at (2.5,-0.65) {基础设施接入};
\node[draw,circle] (d) at (0,-2) {甲};
\node[draw,circle] (e) at (2,-2) {乙};
\node[draw,circle] (f) at (4,-2) {丙};
\node[draw,rounded corners] (g) at (6,-2) {网关};
\draw[<->] (d)--(e); \draw[<->] (e)--(f); \draw[<->] (f)--(g);
\draw[<->] (e) to[bend left=28] (g);
\node at (3,-2.75) {网状转发：先到邻居，再寻找出口};
\end{tikzpicture}
\caption{连接方式的原创示意。线表示可用连接，不表示所有连接同时发送；网状网络中的一条路可能包含多次无线传输。}
\end{figure}

从覆盖尺度看，蓝牙（Bluetooth）常用于个人设备附近的短距离连接，Wi-Fi 常覆盖房间、楼宇或局部室外区域，蜂窝网络（Cellular Network）以基站服务较大的区域，卫星网络（Satellite Network）可以提供广域覆盖。这里比较的是课程中的典型尺度，不给每项技术设定固定不变的距离；天线、功率、环境和具体部署都会改变实际范围。游牧式移动（Nomadic Mobility）强调用户换了位置以后重新接入，例如从教室到宿舍后再使用网络；连续移动中的业务保持，则需要网络在移动期间支持接入变化。

覆盖范围与移动性（Mobility）是两个不同维度。较大的覆盖区通常让移动终端在同一个接入节点下停留更久，但移动支持还依赖网络是否提供合适的管理机制。切换（Handover/Handoff）是终端从一个接入节点转向另一个接入节点的过程。课程比较 Wi-Fi 与蜂窝网络时，既考虑覆盖区大小，也考虑切换管理，而不是单独由信号传播距离推断全部移动能力。至此，网络已经被拆成节点、链路与连接结构；接下来把目光缩到一条链路，研究信息究竟如何穿过空气。

\subsection{从比特到电磁波的完整路线}
比特（Bit）是取值为 0 或 1 的信息单位。空气中传播的不是打印出来的“0”和“1”，而是随时间变化的电磁波（Electromagnetic Wave）。发送设备必须把比特转换成可发射的信号，接收设备再根据信号判断发送端的选择。物理层（Physical Layer, PHY）负责这类实际信号的生成、传播接口和恢复工作。

调制（Modulation）就是让信号的某些特征随信息发生可辨认的变化。数字调制（Digital Modulation）的“数字”说明输入信息是离散比特，不说明实际无线波形只具有两个数值。即使输入只有 0 和 1，天线上产生的信号仍然是连续变化的物理量。接收端进行解调（Demodulation），依据事先约定的信号集合恢复比特。

课程用两个处理步骤解释发射过程。第一步把比特映射成承载信息的波形；第二步把该波形放到选定的无线工作频段。第二步在课件中称作模拟调制（Analog Modulation），其重点是载波上的频率搬移。这里不应把“模拟调制”误读为“用户的信息一定是模拟信息”。同一条发送链路完全可以先处理数字比特，再通过连续波形和射频电路发射出去。

\begin{figure}[htbp]
\centering
\begin{tikzpicture}[>=stealth,every node/.style={font=\small},x=1cm,y=1cm]
\node[draw,align=center] (a) at (0,0) {信息比特\\0101\ldots};
\node[draw,align=center] (b) at (2.3,0) {数字调制\\选择信息波形};
\node[draw,align=center] (c) at (5.1,0) {频率搬移\\进入工作频段};
\node[draw,align=center] (d) at (7.7,0) {发射天线};
\draw[->] (a)--(b);\draw[->] (b)--(c);\draw[->] (c)--(d);
\node[draw,align=center] (e) at (7.7,-1.7) {接收天线};
\node[draw,align=center] (f) at (5.1,-1.7) {搬回基带};
\node[draw,align=center] (g) at (2.3,-1.7) {数字解调\\判断发送信号};
\node[draw,align=center] (h) at (0,-1.7) {恢复比特};
\draw[->] (d)--node[right,align=left] {无线\\传播} (e);
\draw[->] (e)--(f);\draw[->] (f)--(g);\draw[->] (g)--(h);
\end{tikzpicture}
\caption{发送和接收是方向相反的处理链。图暂未加入纠错编码；后文将在比特与调制之间加入它。传播可能损坏信号，所以“恢复比特”并不保证每个比特都正确。}
\end{figure}

\section{频率、频谱与调制：信息怎样进入波形}
\subsection{三个可以改变的波形特征}
先用一个未承载复杂信息的正弦信号（Sinusoidal Signal）描述基本量：
\begin{equation}
s(t)=A\cos(2\pi f t+\phi).
\end{equation}
这里，$t$ 是时间，单位为秒；$A$ 是幅度（Amplitude），表示波形偏离零值的大小；$f$ 是频率（Frequency），单位为赫兹（Hertz, Hz），表示每秒振荡多少次；$\phi$ 是相位（Phase），单位为弧度，表示振荡在一个周期内的起始位置。$s(t)$ 可以理解为与发射信号对应的电压或电场分量，具体单位随所选物理量而定。这个表达式用于解释波形，不把幅度直接当作功率。

周期（Period）$T$ 是完成一次振荡所需的时间，因此
\begin{equation}
T=\frac{1}{f}.
\end{equation}
例如，频率越高，同一小段时间内能看到的振荡次数越多；相位改变则相当于沿时间轴移动波形。两个频率和幅度相同、相位相差 $\pi$ 弧度，即 $180^\circ$ 的正弦波，在任意时刻的取值互为相反数。

发送端可以分别改变这三个量。幅移键控（Amplitude Shift Keying, ASK）以不同幅度表示不同信息；频移键控（Frequency Shift Keying, FSK）以不同频率表示信息；相移键控（Phase Shift Keying, PSK）以不同相位表示信息。在最简单的二进制例子中，FSK 约定频率 $f_0$ 表示比特 0，频率 $f_1$ 表示比特 1。接收端观察一个约定好的时间区间，判断这一段波形更接近哪种频率，再输出对应比特。

关键是双方共享同一张映射表。某种波形本身没有天然的“0”或“1”含义；发送端和接收端必须对波形、观察时间以及映射规则达成一致。调制的过程不是把数字写在波上，而是在允许的信号集合中作出选择。

\subsection{载波频率不等于带宽，也不等于数据速率}
载波（Carrier）是用于承载并搬移信息的高频振荡，其频率称为载波频率（Carrier Frequency）$f_c$。基带（Baseband）是频率搬移之前、围绕较低频率表示信息的信号范围。把基带信号搬到 $f_c$ 附近后，它占用的是一段频率范围，而不仅是一个没有宽度的频率点。

频谱（Spectrum）描述信号包含哪些频率成分，以及这些成分的相对强弱。频带（Frequency Band）指一段频率区间。带宽（Bandwidth）描述这段区间的宽度；在用上下边界近似描述信号占用频带时，
\begin{equation}
B=f_H-f_L.
\end{equation}
其中，$f_L$ 为下边界频率，$f_H$ 为上边界频率，$B$ 为带宽，三者均以 Hz 为单位。实际信号的频带边界需要结合具体占用标准确定；这里采用理想化区间，是为了分清位置和宽度。

假设一个教学信道（Channel）从 $2.400$ GHz 延伸至 $2.420$ GHz，它的带宽是 $0.020$ GHz，即 $20$ MHz，而中间位置是 $2.410$ GHz。这个数值例子只说明单位换算和区间含义，不指定某个实际 Wi-Fi 信道的完整参数。把整个区间移到更高频率，若上下边界一起移动相同距离，带宽仍然不变。

\begin{figure}[htbp]
\centering
\begin{tikzpicture}[>=stealth,every node/.style={font=\small}]
\draw[->] (0,0)--(10,0) node[right] {频率};
\fill[blue!12] (2,0) rectangle (7,1.2);
\draw[thick] (2,0)--(2,1.2)--(7,1.2)--(7,0);
\draw[dashed] (4.5,0)--(4.5,1.6);
\node[below] at (2,0) {$f_L$};\node[below] at (7,0) {$f_H$};
\node[above] at (4.5,1.6) {$f_c$：频带所在位置};
\draw[<->] (2,-0.8)--node[below] {$B=f_H-f_L$：频带宽度} (7,-0.8);
\end{tikzpicture}
\caption{中心频率与带宽的区别。方框是用于理解占用区间的示意，不表示所有真实信号的频谱均为矩形。}
\end{figure}

比特率（Bit Rate）$R_b$ 则表示每秒发送多少比特，单位是 bit/s。GHz 与 Mbit/s 度量的对象不同：前者是振荡或频谱位置，后者是信息传送速度。不能看到“工作在 5 GHz”就推断“速率为 5 Gbit/s”。调制方式、可用带宽、编码开销、信道条件及接入等待共同影响最终的数据传输表现。

频谱资源需要协调使用。授权频谱（Licensed Spectrum）允许特定使用者在规定条件下使用某些频带；免许可频谱（Unlicensed Spectrum）也受使用规则约束，并不意味着没有发射功率限制或没有干扰。课程借频谱许可的成本强调：频率资源有限，网络设计不仅要让一对设备通信，还要让许多用户有效共享资源。

\subsection{从一个比特扩展到一个符号}
符号（Symbol）是在一个约定的符号时间内选择并发送的一种信号状态。若每个符号只在两个信号之间选择，它承载一个比特。若每次能可靠区分八种信号，就可以把三个比特作为一组，因为三个比特具有 $2^3=8$ 种组合。

一般地，对于具有 $M=2^m$ 个可区分状态、采用固定长度映射的调制，
\begin{equation}
m=\log_2 M,\qquad R_b=mR_s=\frac{m}{T_s}.
\end{equation}
这里，$M$ 是信号状态数；$m$ 是每符号承载的比特数；$R_s$ 是符号率（Symbol Rate），单位为符号每秒，也称波特（Baud）；$T_s=1/R_s$ 是符号持续时间，单位为秒；$R_b$ 是映射到这些符号上的比特率。这个式子尚未扣除编码冗余和协议开销，因此不能直接把它当作应用有效吞吐量。

例如，同样每秒发送一百万个符号，二进制调制每个符号带一个比特，对应 $1$ Mbit/s；16 个状态的调制每个符号带四个比特，对应 $4$ Mbit/s。提高速率的原因不是每个波形振荡了四倍次数，而是每次选择能够表达的可能性增加了。与之相伴的困难是：接收端必须更精细地区分这些可能性。

信号星座（Signal Constellation）把允许发送的信号状态画成若干点。点之间距离越近，受扰动后越容易混淆。课程先比较相位相差 $180^\circ$ 的两个状态，再把相位状态扩展到八个。若所有点都放在相同半径的圆上，状态数增加就意味着相邻相位间隔减小。

\subsection{QAM 为何更高效，也更依赖信道质量}
正交幅度调制（Quadrature Amplitude Modulation, QAM）同时利用幅度和相位的不同组合。常用星座图采用同相分量（In-phase Component, I）与正交分量（Quadrature Component, Q）作为两个坐标；这两个分量对应相差 $90^\circ$ 的参考振荡。点到原点的距离体现幅度，点相对坐标轴的角度体现相位。

16QAM 有 16 个允许状态，因此每符号承载 4 个比特；256QAM 有 256 个允许状态，因此每符号承载 8 个比特。这里的 16 和 256 是信号集合的大小，不是频率数，也不是发射天线数。

\begin{figure}[htbp]
\centering
\begin{tikzpicture}[>=stealth,scale=0.72,every node/.style={font=\small}]
\draw[->] (-4.3,0)--(4.6,0) node[right] {I};
\draw[->] (0,-4)--(0,4.4) node[above] {Q};
\foreach \x in {-3,-1,1,3}{\foreach \y in {-3,-1,1,3}{\fill[blue!70!black] (\x,\y) circle (0.09);}}
\draw[dashed,gray] (2,-3.6)--(2,3.6);
\draw[->,thick,red!70!black] (1,1)--(2.5,1.25);
\fill[red!70!black] (2.5,1.25) circle(0.08);
\node[align=left,anchor=west] at (4.4,2) {蓝点：允许发送的状态\\红点：一次受扰动的观测\\虚线：两列之间的示意边界};
\end{tikzpicture}
\caption{原创 16QAM 示意，坐标为相对值。发送点 $(1,1)$ 受到足够大的扰动后，观测点可能落到相邻状态一侧；图不指定具体比特标记或接收算法。}
\end{figure}

接收信号会受到衰减（Attenuation）、噪声（Noise）和干扰（Interference）的影响。衰减让信号变弱；噪声来自不可避免的随机扰动；干扰来自其他发射端的信号。接收端实际观察到的点不一定正好落在发送星座点上，只能根据观测结果判断最可能发送了哪个状态。如果在相同平均发射功率下把更多点放入有限的星座范围，点间间隔往往缩小，错误判断的风险随之上升。

所以，更大的星座不是无条件的改进。信道好时，可以使用更大的星座，提高每符号的信息量；信道差时，较少的状态更容易区分。所谓稳健性（Robustness）与效率的权衡，就是在当前信号质量下决定“每次选择表达多少信息”，而不是孤立地追求最大的调制阶数（Modulation Order）。

\section{无线传播：为什么收到的信号会变弱、变形}
\subsection{频率、波长与天线尺寸}
波长（Wavelength）$\lambda$ 是同一时刻波形沿传播方向重复一次的空间距离。对于传播速度为 $c$ 的电磁波，
\begin{equation}
\lambda=\frac{c}{f}.
\end{equation}
这里，$\lambda$ 的单位为米，$f$ 的单位为 Hz。在真空中 $c$ 约为 $3\times10^8$ m/s；对本节空气中传播的估算，使用这个近似值。频率越高，波长越短。例如 $2.4$ GHz 对应约 $0.125$ m，即 $12.5$ cm；这不是说所有 $2.4$ GHz 天线都必须长 $12.5$ cm，而是说明天线设计的尺寸尺度与波长相关。

天线（Antenna）把电路中的信号与空间中的电磁波联系起来。很多天线的特征尺寸是波长的某个比例，因此高频、短波长有利于在小型设备中布置天线，甚至布置多个天线。不过，高频传播也可能带来更大的损耗。课程在这一处强调的是两个方向相反的因素：小天线有利于移动设备，而更容易衰减的信号又不利于远距离覆盖。

\subsection{功率、能量与分贝要分别理解}
功率（Power）表示每单位时间传递的能量，单位为瓦特（Watt, W）；能量（Energy）的单位为焦耳（Joule, J）。如果一段持续时间 $\Delta t$ 内功率近似恒定为 $P$，则
\begin{equation}
E=P\Delta t.
\end{equation}
这里 $E$ 是该段信号的能量，$P$ 以 W 为单位，$\Delta t$ 以 s 为单位。若比特率为 $R_b$，每比特时间 $T_b=1/R_b$，并在相同平均功率意义下记账，则每比特能量（Energy per Bit）为
\begin{equation}
E_b=PT_b=\frac{P}{R_b}.
\end{equation}
因此，在功率固定时，提高每秒发送的比特数，会减少分摊到每个比特的能量。使用这个式子时，还必须说清 $R_b$ 指编码后的发送比特率还是用户信息比特率；两者不同，所定义的每比特能量也不同。本节先以进入调制器的比特率为准。

分贝（Decibel, dB）用于表示比值。两个正功率 $P_2$ 与 $P_1$ 的比值写成
\begin{equation}
L_{\mathrm{dB}}=10\log_{10}\!\left(\frac{P_2}{P_1}\right).
\end{equation}
如果 $P_2$ 是 $P_1$ 的十倍，结果为 $10$ dB；如果只有十分之一，结果为 $-10$ dB。这个量只表达两个功率的相对关系，不能脱离参考值充当绝对功率。表示路径损耗时通常将大功率放在分子，使损耗为正；表示接收相对发射的变化时则可能得到负数。计算前先确定比值的方向，比机械记忆符号更重要。

以 $1$ mW 为参考的分贝毫瓦（Decibel-milliwatt, dBm）用于表示绝对功率：
\begin{equation}
P_{\mathrm{dBm}}=10\log_{10}\!\left(\frac{P}{1\,\mathrm{mW}}\right).
\end{equation}
例如 $1$ mW 是 $0$ dBm，$10$ mW 是 $10$ dBm，$0.1$ mW 是 $-10$ dBm。负的 dBm 不表示负能量，只说明功率小于参考功率。还要注意两个 dBm 功率不能直接相加作为总功率；应先换回线性功率相加，再按需要转回 dBm。

\subsection{自由空间模型说明了什么}
自由空间传播（Free-space Propagation）先排除建筑、树木和反射路径，假设发射与接收之间存在无遮挡传播。各向同性天线（Isotropic Antenna）是假想的参考天线，它向所有空间方向均匀辐射。定向天线（Directional Antenna）则把更多能量集中到某些方向。天线增益（Antenna Gain）比较某方向与参考天线的辐射或接收能力，不是凭空产生能量。

在自由空间、远场（Far Field）、天线方向与极化匹配且其他损耗忽略的条件下，课件的接收功率关系可写成
\begin{equation}
P_r=P_tG_tG_r\left(\frac{\lambda}{4\pi d}\right)^2.
\end{equation}
其中，$P_t$ 是送入发射天线的功率，$P_r$ 是接收到的功率，二者用相同功率单位；$G_t$ 和 $G_r$ 分别为收发天线在相应方向上的线性增益，不能直接代入 dB 数值；$d$ 是传播距离，单位为米；$\lambda$ 是波长，单位也为米。远场条件要求不能把天线极近距离区域直接套入这个模型。

把不含天线增益的自由空间路径损耗（Free-space Path Loss）记为 $L_{\mathrm{fs}}$，有
\begin{equation}
L_{\mathrm{fs}}=\left(\frac{4\pi d}{\lambda}\right)^2
=\left(\frac{4\pi df}{c}\right)^2,
\qquad P_r=\frac{P_tG_tG_r}{L_{\mathrm{fs}}}.
\end{equation}
$L_{\mathrm{fs}}$ 是无量纲线性比值。固定频率、增益和发射功率时，距离增加为原来的两倍，接收功率降至四分之一；固定距离、增益和发射功率时，频率增加为两倍，在这一模型下接收功率也降至四分之一。两种变化都相当于约 $6$ dB 的额外路径损耗。

“衰减与频率平方成正比”必须连同这些固定条件一起理解。若换频率的同时也改变了天线增益、天线尺寸或传播环境，不能仅凭 $f^2$ 推断实际接收功率。这里的公式提供一个清晰基准，真实链路还要面对其他传播因素。

\subsection{从覆盖范围到空间复用}
高频带来的较小覆盖区，对单条远距离链路可能不利，对一个密集网络却可能有另一面。空间复用（Spatial Reuse）是让空间上相互影响足够小的连接重复使用相同频率资源。网络密集化（Network Densification）通过增加接入节点、缩小单节点服务区域，为更多局部连接创造复用机会。

课程用多个基站重复使用同一组频带来说明这一点。设可用频带数为 $L$，每个频带在给定条件下可服务 $U$ 个用户，那么一个基站的示意容量是 $LU$。如果有 $K$ 个基站能够在互不造成不可接受干扰的条件下重复使用这些频带，理想化总服务量可达到 $KLU$。这里 $L$、$U$、$K$ 均为计数，$LU$ 和 $KLU$ 只是课堂容量示意，不是实际网络的保证值。

增加基站并不自动把容量乘上基站数。相邻同频发射会相互干扰，覆盖、发射功率和资源使用都需要协调。这解释了为什么课程在讨论“高频损耗较大”之后，紧接着讨论“高频有利于更小覆盖和更多复用”：前者评价单条链路，后者评价整个区域内的总服务能力，两者的评价对象不同。

\subsection{障碍物、多径与衰落}
真实传播环境会吸收、散射或改变电磁波的传播方向。大气吸收（Atmospheric Absorption）与雨雾等因素能增加损耗；折射（Refraction）改变传播方向；反射（Reflection）使信号在表面改变行进方向；衍射（Diffraction）使障碍物边缘附近也可能出现传播分量。课程用建筑、地面、山体和树叶说明：传播条件不仅与两端距离有关，还与路径上的环境有关。

大气造成的附加损耗也与频率有关，并不在所有频率上平滑、等量地增加。课件中的雨、雾与水蒸气曲线用于说明，不同环境因素具有不同的频率响应，某些频率附近可能出现明显吸收。因此，选择工作频率时，除了自由空间损耗，还必须考虑实际路径所穿过的介质。课程举树叶季节变化的例子，也是在说明：两端设备不动，不代表传播环境恒定。

地球曲率还会限制视距传播（Line-of-Sight Propagation）。在课件采用的等效地球半径近似中，一根高度为 $h$ 的天线到无线电地平线（Radio Horizon）的距离近似为
\begin{equation}
d_{\max}\simeq\sqrt{2KR_T h}.
\end{equation}
这里，$R_T$ 是地球半径，课件取约 $6370$ km；$K$ 是描述大气折射影响的无量纲等效地球半径系数；$h$ 是天线相对当地地面的高度；$d_{\max}$ 为单端天线到地平线的近似距离。$R_T$ 与 $h$ 必须采用相同长度单位，结果才具有同一单位。此近似使用 $h\ll KR_T$，并不包含山体遮挡或接收灵敏度的限制。它描述几何与折射意义上的视距，不保证在该距离内功率一定足够通信。两端天线都有明显高度时，还需要考虑两端各自的地平线范围，不能把单端式直接当作所有链路的最大距离。

多径传播（Multipath Propagation）意味着接收端收到了同一个发射信号沿不同路径到达的多个副本。不同路径具有不同长度，因此副本到达的时间、幅度和相位可能不同。它们在接收端叠加，既可能增强，也可能削弱总信号。

\begin{figure}[htbp]
\centering
\begin{tikzpicture}[>=stealth,every node/.style={font=\small}]
\node[draw,rounded corners] (tx) at (0,0) {发送端};
\node[draw,rounded corners] (rx) at (8,0) {接收端};
\draw[->,thick] (tx)--node[below] {直接路径 $d_0$} (rx);
\draw[thick] (2.7,2)--(5.3,2);
\node[above] at (4,2) {反射表面};
\draw[->,blue!70!black] (tx)--(4,2)--node[right] {较长路径 $d_1$} (rx);
\node[align=center] at (4,-1.2) {同一信号的副本先后到达\\接收端把它们叠加，而不是分别“挑选”一条};
\end{tikzpicture}
\caption{原创两径示意。真实环境可存在更多路径；图中反射只用于解释不同路径长度与接收叠加。}
\end{figure}

若只考虑路径长度造成的差异，路径差 $\Delta d=d_1-d_0$ 对应时延差 $\Delta\tau=\Delta d/c$，以及载波相位差
\begin{equation}
\Delta\phi=2\pi f\Delta\tau=\frac{2\pi\Delta d}{\lambda}.
\end{equation}
这里 $d_0,d_1,\Delta d,\lambda$ 均以米为单位，$\Delta\tau$ 以秒为单位，$\Delta\phi$ 以弧度为单位；相位按 $2\pi$ 周期重复。真实反射还可能引入额外相位变化，因此该式只提取“多走了一段路”的影响，不能把路径差本身直接当成以弧度计的相位差。

在两个副本幅度相同、频率相同的简化情形下，同相到达会相长叠加（Constructive Interference）；相位相差 $\pi$ 会相消叠加（Destructive Interference）。只有幅度也相同的理想情形才会完全抵消。一般情形只是部分增强或部分减弱。接收位置稍有改变，路径差就可能改变，所以“离接入点更近”不保证每一个局部位置的信号都更强。

衰落（Fading）描述接收信号随位置或时间变化的起伏。阴影衰落（Shadowing）与建筑、地形等遮挡有关，常表现为较缓慢的变化；多径引起的快速起伏在课程中称为快衰落（Fast Fading）。慢衰落（Slow Fading）与快衰落是相对于观察或通信过程的时间尺度而言的。静止设备也可能因周围人或障碍物移动而经历变化，移动设备则更容易不断进入不同的传播组合。课程此处建立物理直觉，不要求进一步计算详细的统计衰落模型。

\subsection{用 SNR 和 SINR 把传播联系到解调}
接收功率大并不必然意味着容易解调：如果噪声和干扰也很强，有用信号仍然难以分辨。信噪比（Signal-to-Noise Ratio, SNR）比较有用接收功率与噪声功率：
\begin{equation}
\mathrm{SNR}=\frac{P_r}{N}=\frac{P_r}{N_0B}.
\end{equation}
这里 $P_r$ 是经历传播损耗后到达接收端的有用功率；$N$ 是接收带宽内总噪声功率；$N_0$ 是噪声功率谱密度（Noise Power Spectral Density），单位为 W/Hz；$B$ 是所考虑的接收带宽，单位为 Hz。式子 $N=N_0B$ 假设带内噪声谱密度近似平坦，并采用课件中与该带宽相配的谱密度约定。所有功率都应以线性单位代入。

如果同时存在其他发射端，就使用信干噪比（Signal-to-Interference-plus-Noise Ratio, SINR）：
\begin{equation}
\mathrm{SINR}=\frac{P_r}{I+N_0B}.
\end{equation}
其中 $I$ 是接收端看到的总干扰功率，也就是各干扰信号经历其各自传播损耗之后的功率总和，不是邻居设备标称的发射功率总和。分母里的噪声与干扰必须使用相同线性功率单位相加。

例如，在统一的任意线性功率单位下，$P_r=100$、$N=1$ 时，SNR 为 $100$，即 $20$ dB；若新增干扰 $I=9$，SINR 变为 $100/(9+1)=10$，即 $10$ dB。有用信号并未变弱，接收质量却显著下降。这正是商场中多个接入点以及共用频段的不同无线技术可能互相影响的原因。

误比特率（Bit Error Rate, BER）描述接收比特判错的比例或相应概率。对于确定的调制、编码和信道条件，较差的 SNR 或 SINR 通常意味着更高的误比特率。但不存在适用于所有调制方式的单一“正确接收阈值”；高阶星座需要更好的区分条件。到这里，传播与调制形成闭环：传播决定观测信号有多可靠，调制决定在该可靠程度下能承载多少信息。

\section{错误控制：增加冗余，还是重新发送}
\subsection{检错与纠错解决不同的问题}
错误检测（Error Detection）判断数据中是否存在可检测的错误；错误纠正（Error Correction）进一步恢复原来的数据。发现“答案不对”不等于知道“应该改哪一位”。当接收端只有受损信息本身时，通常缺少辨别依据，因此发送端需要加入冗余（Redundancy）。冗余不是独立的新用户信息，而是由原有信息按规则生成的附加约束。

前向纠错（Forward Error Correction, FEC）通过信道编码（Channel Coding）预先加入冗余，使接收端在能力范围内直接修复错误。自动重传请求（Automatic Repeat reQuest, ARQ）则主要用检错结果决定是否重新传送。两种思路都付出代价：FEC 要传额外比特，ARQ 可能需要额外发送和等待。实际机制可以组合使用，但课程在这里先分别理解。

\subsection{分组码与码率：冗余究竟占多少}
分组码（Block Code）把 $k$ 个信息比特编码成 $n$ 个发送比特，形成码字（Codeword）。在课件采用的信息位加冗余位的表示中，其中 $n-k$ 个是冗余比特。码率（Code Rate）定义为
\begin{equation}
r_c=\frac{k}{n},\qquad 0<r_c\leq 1.
\end{equation}
这里 $k$ 和 $n$ 是每个块的比特数，$r_c$ 是无量纲比例。码率越低，单位发送比特中用于承载新信息的比例越小，但更多冗余可为保护信息提供空间。它不意味着“码率低就一定在所有条件下更好”：实际检测和纠正能力还取决于码的具体结构。

若编码后的发送比特率为 $R_{\mathrm{coded}}$，在忽略其他开销时，信息比特率为
\begin{equation}
R_{\mathrm{info}}=r_cR_{\mathrm{coded}}=r_cmR_s.
\end{equation}
$R_{\mathrm{info}}$ 和 $R_{\mathrm{coded}}$ 均以 bit/s 为单位；$m$ 是前文定义的每调制符号比特数，$R_s$ 为符号率。这个式子把调制与编码接到一起：大星座提高每符号可发送的比特数，编码冗余则决定其中有多少属于用户信息。协议首部、确认与重传仍未计入，所以应用真正得到的吞吐量还可能更低。

\subsection{重复码：最直观的纠错过程}
重复码（Repetition Code）把同一信息发送多次。三重复码把信息位 0 编为 000，把信息位 1 编为 111，因此 $k=1$、$n=3$、$r_c=1/3$。接收端对这三个位置作多数判决（Majority Decision）：至少两个位置为 0 就判为 0，至少两个位置为 1 就判为 1。

例如，发送 111，接收成 101，虽然中间一位出错，仍有两票支持 1，因此能够恢复原信息。若接收成 001，则已有两位出错，多数判决会错误地输出 0。因此，这个例子每组三位能保证纠正一个翻转错误，不能保证纠正两个错误。“重复三次”提升可靠性的代价，是用三个发送比特承载一个信息比特。

重复码的价值主要在于说明编码如何利用附加结构。接收端知道合法发送结果只能是 000 或 111，才有依据把 101 解释成 111 的受损版本；若任意三位串都合法，就无法据此纠正。更复杂的编码同样依靠合法信号集合受到约束这一原理，并不要求简单地逐位重复。

\subsection{奇偶校验：为什么能发现，却未必能修复}
奇偶校验（Parity Check）给一组信息比特增加一个校验位（Parity Bit）。偶校验要求包含校验位在内的整个码字中，1 的数量为偶数。比如信息串 1011 含有三个 1，所以附加校验位 1，得到 10111，整个码字有四个 1。

如果传输中恰好一位翻转，1 的数量的奇偶性一定变化，接收端就能发现不满足约束。但它并不知道哪一位翻转，因为每个位置的单比特错误都会造成相同的“奇偶性不对”。因此，单个奇偶位可以检测任意奇数个翻转错误，却不能仅靠这个结果定位并纠正一个错误；偶数个翻转错误可能保持原来的奇偶性，从而漏检。

校验位本身同样是通过信道发送的，也可能出错。课堂中为了展示比较步骤，有时先假设收到的校验位正确；完整的偶校验规则则是对收到的整个码字检查奇偶性。无论翻转发生在信息位还是校验位，一个翻转都会破坏整个码字的偶校验约束。

二维奇偶校验（Two-dimensional Parity Check）把信息位排列成矩阵，分别给每行和每列增加校验。若只有一个信息位翻转，就会恰好出现一行不通过、一列不通过，两者交点定位出错位置。

\begin{center}
\begin{tabular}{c|ccc|c}
 & 第1列 & 第2列 & 第3列 & 行偶校验位\\\hline
第1行 & 1 & 0 & 1 & 0\\
第2行 & 0 & 1 & 1 & 0\\
第3行 & 1 & 1 & 0 & 0\\\hline
列偶校验位 & 0 & 0 & 0 & \\
\end{tabular}
\end{center}
在这个原创例子中，假设第 2 行第 2 列由 1 翻成 0，且校验位都正确，那么第 2 行与第 2 列的 1 的数量都变为奇数，其余行列仍通过。接收端找到交点并翻转该位，完成单个信息位的纠正。这个简化示例未在右下角加入校验位的交叉保护，因此不据此声称能处理任意校验位错误，也不声称能纠正任意多个错误。

卷积码（Convolutional Code）是课件提到的另一类结构。信息比特依次进入移位寄存器（Shift Register），输出由当前输入及保留的历史状态共同决定。课件以每次输入 $k$ 位、输出 $n$ 位描述其码率 $k/n$。此处掌握它与独立块处理的区别即可；本段课堂没有展开具体译码算法，因此不引入额外计算程序。

\subsection{从比特串到数据包与 CRC}
数据包（Packet）把一组比特作为一个完整传送单位。首部（Header）携带地址等控制信息，有效载荷（Payload）携带需要交付的用户内容。信件的信封与正文可以帮助理解这种分工：地址使网络知道如何处理这份数据，而正文才是用户主要想交付的内容。

循环冗余校验（Cyclic Redundancy Check, CRC）根据受保护的数据计算检错值。课堂把它描述为有效载荷的简短“摘要”：发送端根据原载荷生成 CRC，并把它与数据一起发出；接收端根据实际收到的载荷重新计算 CRC，然后比较结果。这里的摘要只是功能类比，不表示它能像正文一样概括语义。

课程为解释流程，把 CRC 放在首部，并假定首部能够正确接收。这个安排是教学模型，不意味着所有协议的 CRC 都位于首部，也不意味着真实首部永远不会出错。理解这一前提后，接收过程就可以按次序读清楚：先提取发送端提供的 CRC，再接收载荷，再用相同规则计算本地 CRC，最后比较两个值。

若两个值不同，说明数据没有通过检查，不能按成功接收处理。若两个值相同，接收端按协议视为检查通过。但有限长度的检错值不能为任意长度数据提供绝对无差错证明，仍可能存在未被检测的错误模式。因此，教材将课堂口语中的“正确收到”精确写为“通过检错，按成功处理”。课程没有进行 CRC 多项式除法，本节也不加入这类运算。

\subsection{停止等待 ARQ：发送、确认、超时、重传}
停止等待（Stop-and-Wait）是本段课程实际讲解的 ARQ 方法。发送端一次发送一个数据包，然后等待接收端的确认（Acknowledgment, ACK）。ACK 是一个控制包，表达前一个包已通过接收端检查。如果等待期间得到了有效 ACK，发送端才继续发送下一个新包；如果等到超时（Timeout）仍没有有效 ACK，就重传（Retransmission）原来的包。

完整过程可以按两个参与者的动作展开。发送端先保留待发送数据包的副本，为载荷生成 CRC，发送包并启动计时。接收端解调后计算本地 CRC；在课程采用的流程中，校验通过便返回 ACK，校验失败便不返回 ACK。发送端在计时期限内收到 ACK，结束当前包的等待并处理下一个包；如果计时到期，则再次发送同一个包，而不是直接跳过它发送新包。

\begin{figure}[htbp]
\centering
\begin{tikzpicture}[>=stealth,every node/.style={font=\small}]
\node at (0,0.4) {发送端};\node at (7,0.4) {接收端};
\draw[->] (0,0)--(0,-6.5);\draw[->] (7,0)--(7,-6.5);
\draw[->,thick] (0,-0.4)--node[above,sloped] {包 A：首次发送} (7,-1.2);
\node[right,align=left] at (7,-1.3) {CRC 失败\\不发 ACK};
\draw[decorate,decoration={brace,amplitude=5pt}] (-0.3,-2.5)--(-0.3,-0.4);
\node[left,align=center] at (-0.65,-1.5) {等待\\超时};
\draw[->,thick] (0,-2.7)--node[above,sloped] {包 A：重传同一包} (7,-3.5);
\node[right] at (7,-3.5) {CRC 通过};
\draw[->,blue!70!black] (7,-3.8)--node[above,sloped] {ACK(A)} (0,-4.6);
\draw[->,thick] (0,-5.0)--node[above,sloped] {包 B：收到确认后才发送} (7,-5.8);
\node[left] at (0,-6.4) {时间向下};
\end{tikzpicture}
\caption{原创停止等待时序图。斜线表示传送需要时间，不按真实比例绘制。第一次失败后，发送端因没有收到确认而超时重传。}
\end{figure}

“自动重传请求”不要求接收端一定发出一个写着“请重传”的消息。在这里，不返回 ACK 与发送端的超时机制结合，就触发了重传。另一方面，没有收到 ACK 也不能严格证明数据包一定没到达：确认本身可能丢失。因此，课件还指出包需要序列号（Sequence Number），以区分新包和重复包。序列号让接收端有依据识别重传，而不把重复收到的同一份内容当成两份新内容。课堂没有展开其他 ARQ 协议，本节也停留在停止等待机制。

发送端通常设有最大重试次数。到达上限仍未完成确认，就丢弃当前包或向上层报告失败，不会无穷次占用链路。这也说明 ARQ 不是“保证所有包必定到达”的魔法：它是在时间、能量和可靠性之间作取舍。链路条件很差时，重传次数增加，不仅消耗发射能量，也拉长交付时间。

停止等待名称中的“等待”揭示了它的主要效率限制：在前一个包得到确认之前，发送端不发送下一个新包。一个成功周期包括数据包发送、信号传播、接收处理以及 ACK 返回等过程。即使物理层瞬间能达到较高比特率，应用仍然可能因为等待和重传而获得较低吞吐量。本段课程重点是正确读懂这个过程，不要求推导额外的停止等待效率公式。

\section{从通信基础走向介质访问控制}
经过前面各节，一条无线传输链路已经具有完整含义：用户数据先形成比特，必要时加入冗余，经调制选出波形，再搬到工作频段，由天线发射；信号经历传播损耗、多径、噪声和干扰后，接收端解调、检错或纠错，必要时通过确认和重传完成可靠交付。

然而，这还没有回答一个网络问题：如果许多设备都想在同一无线信道上发送，谁应当先发？介质访问控制（Medium Access Control, MAC）负责协调对共享通信介质的使用。物理层解决“怎样把数据变成可发送、可接收的信号”，MAC 层进一步解决“什么时候允许一个设备发这个包”。即使某个包已经准备好，设备也可能需要等待合适的发送时机。

课程在 10 月 9 日由此转入 Wi-Fi 与 IEEE 802.11。无线局域网（Wireless Local Area Network, WLAN）提供局部区域内的无线连接；IEEE 802.11 是课程用来解释 Wi-Fi 物理层和 MAC 层工作的标准体系。应用（Application）产生数据之后，数据沿协议栈（Protocol Stack）向下传递，到达 MAC 层，最后交给物理层进行实际发送。接入点收到后，再把数据转入后续网络。

这里要保留两条相互关联、但不能混同的主线。第一条是信号质量：衰减和干扰影响判决，调制与编码决定如何适应这种质量。第二条是共享资源：多个设备竞争同一个介质，接入规则决定谁在何时使用它。后续 Wi-Fi 章节将沿第二条主线展开，但每次讨论碰撞、等待、确认与重传时，都要回到第一条主线，理解它们最终如何影响真实的无线发送。
